\documentclass[sigplan,screen,review,anonymous]{acmart}

\usepackage{booktabs}
\usepackage{enumitem}
\usepackage{microtype}
\usepackage{tabularx}
\usepackage{xspace}

\settopmatter{printfolios=true}
\setcopyright{none}
\renewcommand\footnotetextcopyrightpermission[1]{}
\setlength{\emergencystretch}{1.25em}

\newcommand{\rtdl}{\textsc{RTDL}\xspace}

\title{RTDL: A Restricted-Python DSL for Repurposed RT}
\acmConference[CGO '27]{IEEE/ACM International Symposium on Code Generation
and Optimization}{March 20--24, 2027}{Salt Lake City, Utah, USA}
\acmYear{2027}
\copyrightyear{2027}

\begin{document}

\begin{abstract}
Ray-tracing (RT) hardware increasingly executes non-rendering workloads in
which a geometric hit denotes an application candidate, record, neighbor, or
aggregate contribution. Such meanings impose obligations on coverage,
identity, repeated delivery, termination, and reduction. These obligations
cross source callbacks, data and geometry interpretation, physical traversal,
and result return; local stage legality and machine types do not capture the
whole relation.

We present \rtdl, a restricted-Python domain-specific language and compiler for
bounded repurposed RT computations. Authors declare immutable records and typed
role functions. The compiler constructs canonical Callback IR, checks role
effects, resources, semantic and physical contracts, and supported result
routes, then lowers through a deterministic ABI and generated Numba leaves into
trusted topology-specific CUDA/OptiX wrappers. Prepared execution binds the
admitted program, target, and runtime identities. For fixed families, a result
obligation changes concrete compiler decisions: complete-relation routes reject
unsupported early termination, logical all-hit acceptance generates a payload
update followed by physical intersection rejection, and exact callback IR
selects guarded standard-route specialization.

The implementation exposes two stable constructors and a bounded corpus over
triangle, custom-AABB, curve, and sphere leaves. A historical audit covers nine
project-authored V4 mappings. An initial public-PyOptiX comparison covered
selected stages from three mappings (four operation units); all twelve
complete, first-execute, and prepared observations were slower
(1.080--75.533$\times$). A separately identified remediation then put all ten
registered complete/prepared rows over those units and one long graph scale
within a 1.20$\times$-median/1.35$\times$-maximum engineering envelope; the
8.8--9.3\,s graph row was 1.057$\times$. Separately, two exact standard routes
were 1.077--1.175$\times$ Direct when prepared. \rtdl does not infer profitable
mappings, accept arbitrary Python, synthesize arbitrary topologies, prove
application semantics, or establish easier authoring or broad performance.
\end{abstract}

\begin{CCSXML}
<ccs2012>
<concept>
<concept_id>10011007.10011006.10011041</concept_id>
<concept_desc>Software and its engineering~Compilers</concept_desc>
<concept_significance>500</concept_significance>
</concept>
</ccs2012>
\end{CCSXML}

\ccsdesc[500]{Software and its engineering~Compilers}
\keywords{ray tracing, domain-specific languages, restricted Python, GPU compilation}

\maketitle

\section{Introduction}

Dedicated RT traversal is now used for particle tracking, graph processing,
spatial indexing and joins, neighbor search, clustering, database processing,
and N-body simulation
~\cite{wang2022particle,xiao2025rtgraph,geng2025librts,geng2024rayjoin,zhu2022rtnn,nagarajan2023rtdbscan,shi2025raydb,nagarajan2025rtbarneshut,meneses2026rtcoresurvey}.
These systems show that a bounding-volume hierarchy can accelerate much more
than image synthesis. They also expose a recurring programming problem. A
hardware hit may be only a candidate, one logical row, a neighbor, or one
contribution to a checked aggregate. The corresponding requirements for
filtering, logical identity, repeated delivery, early termination, and
reduction differ by result.

Successful applications already solve concrete instances. RayJoin records all
accepted intersections and physically ignores each one so traversal continues;
RayDB prevents duplicate delivery of a primitive from corrupting aggregation;
RTNN uses different maintenance rules for bounded range and nearest-neighbor
results; LibRTS offers spatial-query handlers for counting and collection
~\cite{geng2024rayjoin,shi2025raydb,zhu2022rtnn,geng2025librts}. These are prior
solutions and motivate, rather than constitute, our contribution. The language
question is how a programmable interface can represent selected obligations so
that they constrain both admissible callbacks and generated execution.

Modern RT systems already provide substantial machinery. OptiX exposes the
callback protocol; PyOptiX and SlangPy bring RT construction to Python; OWL and
Luisa automate pipeline construction; DXR payload qualifiers and Slang
capabilities enforce stage and target constraints
~\cite{parker2010optix,nvidiaPyOptix,slangPyRelease,slangPyDocs,nvidiaOWL,zheng2022luisa,microsoftDXRPAQ,slangCapabilities}.
Rendering languages also connect domain results to renderer-owned execution;
for example, Open Shading Language separates closure construction from
integrator evaluation~\cite{openShadingLanguageDocs}. Our claim is not that
rendering lacks protocols or that prior systems cannot express the same
guarantee. We ask a narrower compiler question:

\begin{quote}
For an author-selected RT mapping, how can a restricted high-level language
express bounded non-rendering computation while making selected result
obligations constrain callbacks, data interpretation, traversal actions, and
result return?
\end{quote}

\rtdl answers for a bounded subset. Authors write immutable records, restricted
role functions, and manifests that name data, resources, and result contracts.
The compiler builds typed Callback IR, checks cross-role and result-route
constraints, generates leaf code and trusted topology wrappers, and binds the
accepted executable to a prepared runtime. The application still owns the RT
formulation, geometric encoding, domain predicates, and semantic oracle.

This paper makes four contributions:
\begin{enumerate}[leftmargin=*,itemsep=1pt,topsep=2pt]
  \item We design a restricted-Python DSL and programming model for bounded
  repurposed RT. It exposes typed records and role-indexed effects while keeping
  raw PTX, SBT, pipeline, and traversal operations compiler-owned.
  \item We develop typed Callback IR and whole-protocol admission that relate
  roles, result obligations, semantic and physical data contracts, resources,
  continuation, and executable identity. This is a fixed-family design, not
  automatic inference or a soundness proof.
  \item We implement a deterministic ABI, generated Numba leaves, trusted
  topology-specific wrappers and exact-IR specialization, plus a prepared
  identity-checked runtime. The lowerers and specializers remain in the TCB.
  \item We reconstruct nine project-authored V4 mappings and evaluate bounded
  mutation, composition, target checking, lifecycle, and performance. An
  adverse selected-stage 3/9 comparison motivates a separately identified
  deployment remediation; ten registered successor rows meet its engineering
  envelope, including one multi-second graph computation. Receipt, startup,
  coverage, generality, and authoring limits remain explicit.
\end{enumerate}

\section{RTDL by Example and Programming Model}

Figure~\ref{fig:dsl-example} shows an abridged version of the supported
built-in-triangle all-hit count callback used by the standard library. The full
program also defines \texttt{make\_ray} and \texttt{miss}; those roles are
omitted only from the figure. The front end parses this text without importing
or executing it as Python.

\begin{figure}[t]
\begin{minipage}{.98\columnwidth}
\footnotesize
\begin{verbatim}
@optix.payload
class CountPayload:
    count: u64

@optix.output
class CountOutput:
    count: u64

@optix.program(
    payload=CountPayload, output=CountOutput,
    attributes=(), max_trace_depth=1,
    max_callable_depth=0)
class TrianglePerRayCount:
    # make_ray and miss omitted here
    @optix.any_hit
    def any_hit(hit: TriangleHit,
                payload: CountPayload) -> AnyHitEffect:
        updated = CountPayload(
            count=payload.count + 1)
        return optix.accept_continue(
            payload=updated)

    @optix.finalize
    def finalize(payload: CountPayload) -> CountOutput:
        return optix.output(value=CountOutput(
            count=payload.count))
\end{verbatim}
\end{minipage}
\caption{Abridged real \rtdl source for a supported all-hit count. The
manifest, ray construction, and miss role are omitted for space; the figure is
not a standalone executable.}
\Description{Restricted Python declares a count payload and output, an any-hit
role that increments the payload and continues, and a finalize role that emits
the count.}
\label{fig:dsl-example}
\end{figure}

The author supplies the callback source, manifest, records, constants, geometry
and data bindings, result contract, and an application oracle. The manifest
also separates \emph{static input}, such as geometry and primitive metadata
prepared once, from \emph{batch input}, such as query rays supplied repeatedly.
The compiler owns parsing, verification, role/effect IR, ABI flattening, leaf
code, trusted wrapper selection, PTX composition, program/SBT binding, status
handling, and prepared lifecycle. Users cannot inject raw traversal calls,
arbitrary PTX, native callback pointers, or self-declared physical authority
through the stable surface.

Compilation follows:
\begin{center}
\small
source + manifest $\rightarrow$ typed Callback IR $\rightarrow$
whole-protocol admission $\rightarrow$ ABI/Numba leaves $\rightarrow$
trusted wrapper/PTX $\rightarrow$ bind/prepare/execute.
\end{center}
The language surface is broader than the two stable closed constructors, but
not every verified callback/topology combination has an executable lowerer.
Unsupported combinations fail closed rather than falling back to arbitrary
Python or handwritten target code.

Table~\ref{tab:prior-obligations} grounds the language problem in three prior
systems that repurpose RT hardware. Each already implements its listed solution.
\rtdl targets the narrower compiler relation between a declared obligation and
its own accepted/generated route; it does not take ownership of the underlying
application predicate or algorithm.

\begin{table*}[t]
  \caption{Recurring result obligations in prior repurposed-RT systems. These
  are motivation and prior solutions, not RTDL inventions.}
  \label{tab:prior-obligations}
  \small
  \begin{tabularx}{\textwidth}{@{}p{.12\textwidth}p{.23\textwidth}p{.29\textwidth}X@{}}
    \toprule
    System & Result obligation & Existing application solution & RTDL's bounded compiler target \\
    \midrule
    RTNN~\cite{zhu2022rtnn} & Range-K may stop at K; nearest-K needs ranked candidates & Intersection filtering plus result-specific bounded storage or priority queue & Represent the selected continuation/result contract; do not infer distance correctness or search coverage \\
    RayJoin~\cite{geng2024rayjoin} & Complete logical intersections despite AABB false positives & Exact predicate, result recording, then physical intersection ignore to continue & Relate logical event acceptance to generated traversal action and bounded publication \\
    RayDB~\cite{shi2025raydb} & A record delivered by multiple rays must not corrupt aggregation & Primitive identity plus first-delivery atomic flag before aggregation & Represent identity/delivery/reducer obligations for supported routes; do not infer SQL bag semantics \\
    \bottomrule
  \end{tabularx}
\end{table*}

\section{IR and Whole-Protocol Checking}

\subsection{Why callback typing is insufficient}

A ray callback runs inside a protocol distributed across host and device code.
Its legal effects depend on its role: an any-hit callback may ignore or
terminate an intersection, whereas a miss callback cannot. Its payload fields
must agree with every producer and consumer. Its program group must match the
acceleration-structure leaf kind. The launch must occur only after module,
pipeline, shader-binding-table, geometry, and output state have been prepared.
Finally, the code checked by the compiler must be the code loaded by the runtime.

These are related to typestate, interface automata, and session-style protocol
reasoning~\cite{strom1986typestate,deAlfaro2001interfaceAutomata,honda2016multiparty},
but the FFI and generated-code boundary adds an identity problem familiar from
linking and foreign interfaces~\cite{furr2005ffi,patterson2017linkingTypes}.
Table~\ref{tab:result-route} makes the additional result-related decisions
concrete. Each row is an implemented fixed-family rule, not a theorem for
arbitrary RT programs.

\begin{table*}[t]
  \caption{Implemented result-route relations and their evidence boundaries.}
  \label{tab:result-route}
  \small
  \begin{tabularx}{\textwidth}{@{}p{.19\textwidth}p{.25\textwidth}p{.34\textwidth}X@{}}
    \toprule
    Observable result obligation & Not decided by stage or machine typing &
    RTDL constraint or transformation & Evidence boundary \\
    \midrule
    Complete bounded relation & Whether locally legal \texttt{TERMINATE} or
    \texttt{IGNORE} fits this route & The family verifier rejects the route
    unless returns are \texttt{ACCEPT\_CONTINUE}; capacity failure publishes no partial
    relation & Source rule and existing CPU checks; not a general enumeration
    theory \\
    Triangle all-hit count & Whether logical acceptance should physically
    accept an intersection & The generator emits a payload update, then calls
    \texttt{optixIgnoreIntersection()}; require single any-hit delivery and
    checked overflow & Generic and specialized lowering source; OptiX
    mechanisms are established, not invented here \\
    Standard count specialization & Whether same-shape \texttt{+1} and
    \texttt{+2} callbacks may share a fixed intrinsic & The generator selects
    the intrinsic only under its exact-IR guard; \texttt{+2} selects the general leaf path & Existing
    source-to-wrapper test replay; no GPU equivalence or general optimization
    theorem \\
    Relation application ID & Whether an attribute denotes physical index or
    application ID despite equal \texttt{u32} width & Admission rejects a
    mismatch between trusted schema and compiled-contract projection & Existing field
    mutation and mock; application intent is not inferred \\
    \bottomrule
  \end{tabularx}
\end{table*}

\subsection{A concrete semantic-ABI failure}

Consider two acceleration-structure primitives at physical positions 0 and 1
whose application identifiers are 10 and 20. Queries 100 and 101 should emit
the pairs $(100,10)$ and $(101,20)$. If an intersection producer instead writes
the physical position into attribute slot zero while the consumer interprets
that slot as an application identifier, the apparent rows are $(100,0)$ and
$(101,1)$. Every value remains a valid \texttt{u32}; width, slot, and row layout
need not expose the semantic substitution. This is an illustrative witness, not
a retained execution of the defective route.

The defect is not a machine-type error but a disagreement over what the bits
mean. In \rtdl, a trusted bounded-relation schema declares that attribute slot
zero denotes an application identifier; the compiled relation contract carries
that row source into a separately derived target projection. Admission compares
the two compiler representations. Existing tests mutate only this fact and
obtain \texttt{CP002\_ATTRIBUTE\_ABI\_OWNERSHIP\_MISMATCH}; with native
initialization overlap disabled, an integrated projection mutation is rejected
before the native-library loader is called. This proves check liveness,
not automatic inference of application meaning or end-to-end execution of the
illustrative defect. Both derivations remain in one compiler TCB, and a wrong
but internally coherent trusted schema can pass.

\subsection{Failure model}

We distinguish four outcomes. \emph{Reject} means admission fails before a
launch. \emph{Status failure} means execution reports failure and output is not
published. \emph{Wrong output} means a worker reports success but an independent
oracle disagrees. \emph{Unverified output} means a value exists without the
required receipt or oracle evidence. The distinction prevents a status gate
from being described as numerical correctness and prevents a generated plan
from being described as an executable lowering.

A focused mock reproduced an unrepaired provider double fault: secondary bind or
close failure can replace the primary exception and lose cleanup retry ownership;
the mock returned no public result. A native-fork probe accepted a public call
through a mock native implementation after bypassing the registered Python
at-fork hook; it did not execute GPU work. Neither defect was observed in
retained successful GPU workers, and inherited owners after such forks remain
unsupported.

\subsection{Typed callback representation}

The source language is restricted Python, not arbitrary Python. Source is
parsed through an AST allowlist and is not imported or executed for discovery.
The front end resolves declared immutable records, frozen constants,
same-unit acyclic helpers, and selected intrinsics, then produces canonical
typed IR. Deserialization reconstructs typed objects and reruns verification;
serialized IR is not itself an authority. Table~\ref{tab:ir} summarizes the
representation.

\begin{table*}[t]
  \caption{Implemented bounded representation. Rejected forms have no fallback
  lowering.}
  \label{tab:ir}
  \small
  \begin{tabularx}{\textwidth}{@{}p{.16\textwidth}p{.40\textwidth}X@{}}
    \toprule
    Form & Accepted subset & Rejected boundary \\
    \midrule
    Types & Booleans; fixed-width integers and floats; 2--4-lane vectors;
    tuples; acyclic immutable records; checked read-only views & Raw pointers,
    mutable containers, opaque objects, and recursive records \\
    Expressions & Typed names and literals, record fields, checked indexing,
    arithmetic, comparisons, bit operations, constructors, selected pure
    intrinsics, and bounded helper calls & Reflection, dynamic dispatch,
    allocation, foreign calls, and unknown attributes \\
    Statements & Typed binding, loop-local type-preserving update, two-sided
    conditionals, statically bounded loops, role-effect return, and helper-value
    return & Dynamic loops, recursion, exceptions, generators, async control,
    and not-all-paths-return \\
    Manifest & Payload/output records, machine attributes, constants, numeric
    policy, resource budget, topology, continuation, and linkage & Self-minted
    physical authority, unreviewed production linkage, and unsupported budgets \\
    Program & Canonical source and IR identities, records, functions, manifest,
    and separately derived effect digest & Unknown fields, malformed schema,
    stale identities, or unverified deserialization \\
    \bottomrule
  \end{tabularx}
\end{table*}

The current resource contract admits at most 32 payload slots and eight
attribute slots, trace depth one, callable depth zero, bounded helper depth, and
statically bounded iteration. The numeric policy requires explicit binary32
behavior and fail-closed handling of prohibited integer overflow and nonfinite
effects. These are language and resource checks; they do not prove that a
chosen ray encoding computes the intended application function.

\subsection{Role-indexed effects}

Seven language roles participate in an admitted unit. \texttt{bounds} runs
during preparation; \texttt{make\_ray} and \texttt{finalize} are language
leaves invoked by a trusted ray-generation wrapper; and \texttt{intersection},
\texttt{any\_hit}, \texttt{closest\_hit}, and \texttt{miss} map to target-stage
behavior. User code cannot invoke raw traversal. Instead, each role returns one
typed effect from a role-specific set: bounds returns an AABB, ray construction
returns a trace request, intersection returns hit or no-hit, hit callbacks
return payload plus continue/ignore/terminate control, miss returns payload,
and finalize returns an output.

Topology makes these requirements relational. Custom AABBs require bounds,
ray construction, intersection, miss, finalize, and at least one hit role.
Built-in triangles forbid bounds and intersection because hardware owns them,
but retain ray construction, miss, finalize, and at least one hit role. An
any-hit role also requires an explicit delivery and continuation contract; the
presence of that declaration does not prove order independence for arbitrary
callback logic.

Role legality is necessary but not sufficient for a fixed result route. The
general any-hit role admits continue, ignore, and terminate effects. The
current complete bounded-relation verifier instead collects all return effects
and requires only logical acceptance with continued traversal; it defines no
filtering or early-termination meaning
for that route. Thus a role-valid effect can be rejected because it conflicts
with the route's output obligation. This is an implemented family restriction,
not a claim that all complete-enumeration designs must reject filtering.

Logical effects also need not map one-to-one to physical intersection actions.
In the triangle all-hit route, the logical continue effect first updates the
payload and then physically ignores the intersection, preserving later
events instead of shrinking traversal's accepted interval. Native geometry
separately requests one any-hit delivery per primitive, and the reduction checks
overflow. These are established OptiX mechanisms; \rtdl's contribution is to
make their required combination part of the trusted implementation of a fixed
result contract, not to invent all-hit traversal.

\subsection{Fail-closed judgments}

Expression checking derives $\Gamma \vdash e : \tau$. Statement checking also
returns the output environment, all-paths-return bit, conservative static work,
and observed effects:
\begin{equation}
  \Gamma;r \vdash s \Rightarrow \langle \Gamma',q,n,\epsilon\rangle .
\end{equation}
For each role $r$, verification enforces its exact signature and
\begin{equation}
  M \vdash F_r : \epsilon
  \quad\text{with}\quad
  \epsilon \subseteq \mathsf{AllowedEffects}(r).
\end{equation}
Names are single-assignment except for type-preserving updates inside static
loops; loop locals cannot escape; both non-returning branches define the same
typed names; and every role/helper path returns. Whole-unit checking validates
schema/source identity, acyclic records and helpers, numeric and resource
budgets, role cardinality, topology, linkage, and required external physical
authority:
\begin{equation}
  \mathcal{G};\mathcal{R} \vdash Q \Downarrow
  \mathsf{Verified}(h_{\mathrm{IR}},h_{\mathrm{effects}}).
\end{equation}
Here $\mathcal{G}$ is an out-of-program geometry authority and $\mathcal{R}$ is
the compiler-owned role topology. The judgment is explanatory notation for
implemented checks, not a mechanized soundness theorem.

Admission is conjunctive. A route is accepted only if its schema is valid, its
roles permit all effects, its semantic and physical ABIs agree, a trusted
lowerer exists for its topology, the generated executable identity is bound,
and the runtime lifecycle can reach the prepared state. Failure of any one
predicate rejects the route; there is no ``best effort'' projection into a
nearby executable.

\subsection{Layered whole-protocol admission}

Table~\ref{tab:admission} identifies where each obligation originates. This
separation matters because target evidence is regenerated within one trusted
implementation; it is not an independent proof of compiler correctness.

\begin{table*}[t]
  \caption{Authorities compared by whole-protocol admission.}
  \label{tab:admission}
  \small
  \begin{tabularx}{\textwidth}{@{}p{.15\textwidth}p{.28\textwidth}p{.28\textwidth}X@{}}
    \toprule
    Obligation & Declaration authority & Target/runtime evidence & Exact limit \\
    \midrule
    Role/effect & Verified role topology and returned effects & Reverified
    compiled callback ABI variants & Common compiler TCB; no arbitrary effect
    inference \\
    Semantic ABI & Trusted family schema and fixed-route nominal row-source
    declaration plus machine types & Compiled relation contract and callback
    ABI & Same compiler TCB; meaning is declared, not inferred; generic
    \texttt{u32} has no nominal meaning \\
    Physical ABI & Typed plan for geometry, buffers, hit channel, result, and
    reducer & Rebuilt target facts and bound storage & Exact-plan admission,
    not arbitrary physical synthesis \\
    Continuation & Completeness/capacity policy and required status-before-use
    order & Status vocabulary, compiled continuation, runtime status gate &
    Bounded supported modes; not a confluence theorem \\
    Identity/lifecycle & Materializer-selected source/PTX/provider identity and
    state machine & Loaded bytes, prepared objects, execution status, output
    digest, and optional receipt & Hash equality is identity coherence, not
    semantic correctness \\
    \bottomrule
  \end{tabularx}
\end{table*}

\section{Code Generation and Runtime}

\subsection{Leaves, wrappers, and specialization}

Verification produces one role-indexed ABI containing flattened inputs,
outputs, effect tags, status tags, and identities. For the general route, the
compiler emits deterministic Python source for each reachable role and helper,
then invokes an isolated Numba compiler process to produce target-specific PTX
device functions. The isolated child receives compiler-generated source rather
than user modules and does not discover the active CUDA device. This leaf step
does not build an OptiX pipeline.

A trusted topology wrapper composes the leaves with compiler-owned raygen,
intersection, hit, and miss mechanics. It checks returned tags before issuing
payload writes, \texttt{optixIgnoreIntersection()}, or termination. Thus the
user effect \texttt{accept\_continue} is a logical action; its physical action
depends on the admitted route. Wrapper source, leaf PTX, ABI, provider symbols,
and physical schema are hashed into the executable identity.

The evaluated standard routes also use trusted specializations guarded by
exact callback IR and reducer identity. For the standard count program in
Figure~\ref{fig:dsl-example}, the specialized entry implements the known
increment-and-continue behavior directly; the general path obtains the same
declared behavior through the leaf ABI and wrapper. Bounded relation similarly
fuses intersection and row emission. These specializations replace generic
leaf calls and per-leaf tag interpretation at selected roles, but retain static
schema, ABI, identity, overflow, capacity, status, and supplied
expected-output checks. An exact-IR guard selects trusted code; it is not a
proof that the specialization preserves source semantics.

Selection depends on program identity, not just role and ABI shape. An existing
compiler test changes the standard increment from one to two. The program
remains front-end legal and has the same role shape, but its IR identity changes
and the generator exits the fixed count intrinsic for the general leaf path.
This source-to-wrapper check establishes one selection boundary; it neither
executes the modified program on a GPU nor proves arbitrary specialization
correctness. The recognizer, wrapper, and specialized lowerer remain in the
trusted computing base (TCB).

\subsection{Plan versus executable lowering}

The canonical planner normalizes the admitted schema and emits a declarative
plan. That plan is intentionally marked non-executable. It does not synthesize
arbitrary CUDA or OptiX callbacks. Instead, an accepted route names a
compiler-owned, topology-specific lowerer and wrapper. The lowerer emits source
and PTX, binds native symbols, constructs the program groups and shader-binding
table, and records their identities. This design supplies a common admission
surface while keeping concrete backend knowledge visible in the TCB.

Shared schema and lifecycle machinery do not make device code generic. The
prospective sphere route in Section~\ref{sec:prospective} required about 2,635
lines of topology-specific code and 28 modified compiler lines. It shows that a
trusted route can reuse three sealed shared files, not automatic topology
lowering.

\subsection{Executable identity and lifecycle}
\label{sec:lifecycle}

The system binds hashes and shape metadata for generated source, PTX, ABI,
native provider symbols, and the public route. The exact app-free native runtime
may be loaded and warmed before route admission, including concurrently with
AOT artifact verification. Admission still gates acceptance of a materialized
or loaded route and its subsequent per-route preparation. The admitted route
lifecycle is:

\begin{center}
\small
static verify $\rightarrow$ materialize/bind $\rightarrow$ prepare
$\rightarrow$ execute/status gate $\rightarrow$ output $\rightarrow$ receipt.
\end{center}

The materialized-program interface validates execution identity, native status,
output digest, and traversal receipt before returning
\texttt{ProtocolExecution}\allowbreak\texttt{Result}. The measured AOT path instead returns
\texttt{RTDLExecutionResult}: it synchronously rejects native and
compact-status failures and any supplied expected-output mismatch, but an
ordinary fast result may omit a digest and detailed receipt. The worker checks
value and digest after prepared return; for first result this is after the
prepared timer but inside the endpoint. It does not rehash disk files per
launch. The preregistered per-execution receipt requirement is unmet: 4,096
timed A executions have only 32 separate post-loop diagnostic executions, one
detailed receipt per A worker.

\begin{figure*}[t]
\centering
\fbox{\begin{minipage}{.95\textwidth}
\small
\textbf{Untrusted program and data}
$\rightarrow$ restricted-Python front end
$\rightarrow$ shared schema, role/effect, ABI, and topology admission
$\rightarrow$ non-executable canonical plan

\medskip
\textbf{Trusted topology route}
$\rightarrow$ topology-specific lowerer and wrapper
$\rightarrow$ generated source/PTX and provider symbols
$\rightarrow$ executable-identity binding
$\rightarrow$ OptiX build, prepare, and launch

\medskip
\textbf{Measured AOT prepared execution}
$\rightarrow$ native status gate
$\rightarrow$ compact status gate
$\rightarrow$ supplied expected-output check, when present
$\rightarrow$ public return (prepared steady timer ends)

\medskip
\textbf{Experiment validation}
$\rightarrow$ returned value
$\rightarrow$ worker output-and-digest oracle
$\rightarrow$ accepted sample or worker success (first-result endpoint ends)

\medskip
\textbf{Separate diagnostic and evidence paths}
$\rightarrow$ post-loop diagnostic execution and one retained detailed receipt
per A worker
$\rightarrow$ offline checker over generated source, PTX, ABI, symbols, and
receipts; it does not import the \rtdl compiler.
\end{minipage}}
\caption{Architecture and trust boundary. Shared admission does not replace the
topology-specific lowerer; that lowerer remains in the TCB.}
\Description{Architecture in which untrusted restricted Python enters shared
admission and a trusted topology-specific lowerer builds the OptiX executable.
Measured AOT status and optional expected-output checks precede return; experiment
validation follows return; separate post-loop diagnostic and offline checker
paths retain bounded evidence.}
\label{fig:architecture}
\end{figure*}

\section{V4 Applications and Case Studies}

RTDL was developed against repurposed-RT applications rather than rendering
shaders alone. Table~\ref{tab:apps} reconstructs nine project-authored V4
mappings to algorithms credited to prior systems. RTDL contributes restricted
expression, checking, generation, and composition of selected parts, not the
application algorithm or geometric mapping. E1 means an exact historical
application-source hash and callback-consumer binding survive. E2 means a
frozen audit recorded the path and structural properties, but an independently
retained historical per-file identity and the original 10.8-MB archive do not;
recovered current entry points do not repair that custody gap. E2 is mapping
evidence, not current runnability.

\begin{table*}[t]
  \caption{Historical V4 mappings. Authors select the RT formulation and own
  the listed application semantics; E1/E2 are custody levels defined in text.}
  \label{tab:apps}
  \scriptsize
  \begin{tabularx}{\textwidth}{@{}p{.12\textwidth}p{.35\textwidth}X@{}}
    \toprule
    Application & Mapping and author-owned semantics & RTDL route and exact boundary \\
    \midrule
    Particle~\cite{wang2022particle} & Tetrahedral closest-face transition; cell encoding, transition, and I/O meaning & E1: built-in triangles, restricted closest-hit update, callback lowering, prepared lifecycle \\
    Triangle~\cite{xiao2025rtgraph} & RT-1A2/RT-2A1 selection; graph orientation and count meaning & E1: all-hit callback, checked capacity/status, segmented reduction; measured scalar is one algorithm stage \\
    RayDB~\cite{shi2025raydb} & Partitioned grouped-I64 aggregate; relation, keys, partitions, and SQL meaning & E2: bounded emission and grouped reduction; sole private-loader exception \\
    LibRTS~\cite{geng2025librts} & AABB point/range containment; predicate, schema, and count/collect meaning & E1: typed custom-AABB relation and capacity/receipt checks; LibRTS is strong prior abstraction \\
    X-HD~\cite{geng2026xhd} & Directed max-of-nearest witness; Hausdorff semantics and output & E2: cell-MBR traversal, checked state, tie-breaking, and reduction \\
    RTNN~\cite{zhu2022rtnn} & Multiround distance-window top-k; metric, $K$, bounds, and ordering & E2: round/refit lifecycle and bounded selection; metric-specific traversal stays trusted \\
    RT-DBSCAN~\cite{nagarajan2023rtdbscan} & Radius graph/components; radius, min-points, and clustering meaning & E2: bounded edges and grouped continuation; not every stage uses RT cores \\
    RayJoin~\cite{geng2024rayjoin} & Six-batch planar overlay; schedule, predicates, and six full result tables & E2: typed producer/reducer route; recovered adapter returns summaries, not the registered full tables \\
    RT-BarnesHut~\cite{nagarajan2025rtbarneshut} & Aggregate hierarchy; tree, acceptance rule, and force law & E2: hierarchy/GAS ABI and frontier checks; tree traversal remains a trusted partner \\
    \bottomrule
  \end{tabularx}
\end{table*}

A frozen AST audit found no selected raw OptiX, CUDA, or PTX assembly tokens in
the nine application files. This supports a structural shift in physical-owner
preparation, not developer-time or LOC savings. Runtime loading passed through
registered V4 interfaces for eight applications; RayDB is the exception, and
trusted C++/CUDA/OptiX partners remain. The stable facade has two constructors,
covers 2/6 OptiX build-input values and 2/4 leaf kinds, while the bounded corpus
covers 4/6 values and four leaf kinds. Only one callback template is presently
authorable, no independent user session was measured, and every new topology
still requires a trusted lowerer.

\section{Evaluation}

We organize existing evidence around five questions. \textbf{RQ1} asks whether
DSL and IR facts change admission, generation, or specialization. \textbf{RQ2}
asks what historical mappings reuse and what remains application-owned.
\textbf{RQ3} asks what a sealed extension reuses and what remains specific to a
topology. \textbf{RQ4} asks what finite target properties an independent checker
can validate. \textbf{RQ5} asks the runtime cost of exact and application paths. Evidence
types remain separate: source traces and component tests are not GPU
equivalence tests, historical mappings are not current runnable examples,
finite mutations are not a soundness proof, and measured routes are not
arbitrary-program lowering.

\subsection{Earlier liveness and residual evidence}

For RQ1, the role-level any-hit set permits continue, ignore, and terminate,
while the complete bounded-relation route accepts only continue returns. The
triangle lowerer supplies the complementary transformation evidence: it records
an accepted event before physically ignoring the intersection to preserve later
events. These are source-traced implemented rules; this paper did not execute a
deliberately terminating complete-relation program on the GPU.

Also for RQ1, we replayed an existing compiler test that changes the standard count
update from \texttt{+1} to \texttt{+2}. Front-end verification still succeeds,
the callback IR identity changes, and wrapper generation switches from the
fixed count intrinsic to the general leaf path. The test exercises parsing, IR
verification, contract/ABI construction, and wrapper generation. It does not
execute either generated path on a GPU or establish semantic preservation for
other optimizations.

For two fixed contracts, 19 single-leaf mutations exercised 19/19 expected
gates: seven role/effect, one semantic-ABI, seven physical-ABI, two continuation,
and two identity mutations. This denominator covers only those contracts and
mutations and is not pooled with later experiments.

With native-initialization overlap disabled, corrupting each target projection
produced the expected sole reason before the native-library loader was called.
For semantic ABI, this is mutation sensitivity, not an
execution of the illustrative wrong-row route. Historical PyOptiX/OWL-style
diagnostics lack independently verifiable execution records here and are not
evidence about those systems' guarantees~\cite{nvidiaPyOptix,nvidiaOWL}. The
19/19 result is neither soundness nor prevalence evidence.

\subsection{Historical application evidence}

For RQ2, the frozen source audit behind Table~\ref{tab:apps} classified all
nine V4 mappings as a structural integration-responsibility shift and found
registered-interface runtime loading in eight. It measured neither authoring
time nor developer productivity. Three exact historical app/callback consumer
identities remain independently recoverable. Current entry points have been
recovered for all nine, but the other six retain only archive-level historical
audit records. We therefore use the portfolio to show bounded mapping and
composition diversity, not historical byte custody, current installability, or
unseen-app success.

A separate historical RTX 4000 Ada authority reported 464 exact,
behaviorally true-OptiX workers and 34 V2-direct/V4 rows. Its row-local median
criterion split 16 pass and 18 fail; 95\% confidence-interval classification
gave 11 clear V4 wins, 10 clear losses, and 13 uncertain. The raw archive is not
part of the current nine-member submission artifact and was not rerun here.
These mixed results are retained to prevent an inference that nine mappings
imply broad performance superiority. They are not a PyOptiX comparison and are
not pooled with the final two-task measurements below.

\subsection{Sealed-snapshot composition exam}
\label{sec:prospective}

For RQ3, before a public randomness pulse, we fixed an author-defined ten-row challenge:
seven built-in four-role and three custom six-role rows. All used supported
primitive families and returned one value per query: six produced counts, and
four produced Booleans by terminating on the first accepted hit. The curve
any-hit-terminate row remained eligible but unselected. The pulse selected
\nolinkurl{builtin_sphere::any_hit_count_continue_u64_per_query}.

Private custody identifies the sealed source. No bytes changed in three frozen
schema, admission, and lifecycle files. The route made two OptiX launches and
passed 12/12 rational-oracle rows, showing reuse of admission, identity, and
lifecycle for this composition. It is not an unbiased application: authors
defined the domain, rows recombined known ingredients, and the selected route
required substantial topology-specific code.

\subsection{Independent finite checker}

For RQ4, the offline checker imports no \rtdl module or compiler projection. It reads
generated source, PTX/ABI metadata, symbols, and receipts. Across three route
groups and four modes, it applied five property classes 20 times using 15 unique
mutations, checking selected target-code patterns and ordering constraints.

This checker is finite, not a verifier for arbitrary device control flow.
During review, an out-of-authority in-memory probe inserted an unconditional
early return and still passed selected effect/status helper checks. The result
therefore supports specialized target-structure validation only. General
control-flow, numerical, and refinement soundness remain open; tools such as
GPUVerify illustrate the substantially stronger proof obligation~\cite{betts2012gpuverify}.

\subsection{Performance methodology}

For RQ5, we measured two exact tasks. \emph{Relation} uses 4,096 objects and 4,096
queries, returns 4,096 canonical u32 pair rows, and performs two true OptiX
traversals. \emph{Triangle} uses 16,384 primitives and 16,384 queries, performs
one true traversal, and returns the checked scalar 65,530.

Five arms isolate different boundaries:
\begin{center}
\small
\begin{tabularx}{.95\columnwidth}{@{}lX@{}}
A & current \rtdl ahead-of-time public path at measured source M; \\
B & idiomatic pinned PyOptiX-compatible baseline; \\
C & strong PyOptiX-compatible baseline with device continuation; \\
D & named Direct CUDA/OptiX reference implementation; \\
E & frozen predecessor control for \rtdl. \\
\end{tabularx}
\end{center}

We ran the same registered matrix on an RTX~4090 (Ada) and RTX~3090 (Ampere).
Each generation has eight blocks, two tasks, five arms, 80 cells, and 128 steady
samples per cell: 160 cells and 20,480 samples total. Ratios are medians of eight
within-block ratios of integer nanoseconds; lower is better, and raw times are
not compared across machines. The steady timer encloses a prepared public action
through return; validation follows. Entry starts before implementation-specific
imports, while post-import starts after them; both end after first-result
validation. PyOptiX initializes CUDA during import, whereas \rtdl does so later,
so these intervals include different lifecycle work.

The reported prepared latencies measure the disclosed exact-standard-IR
triangle and relation specializations, not execution of every role through the
general Numba-leaf ABI.

The registered primary criterion was prepared A/D median $\leq1.20$, with no
maximum-block gate. A/C entry used median $\leq1.20$ and maximum $\leq1.35$, but
the endpoint was revised after an adverse result and both first-result endpoints
remain confounded. We report A/C only diagnostically. Only prepared A/E was a
registered regression gate; A/E first-result values are post hoc. C/B competence
used $\leq1.05$.

\subsection{Prepared execution results}

Table~\ref{tab:prepared} reports the main observation. The public AOT path is
1.077--1.175$\times$ the direct implementation by median. The triangle maximum
block on Ada is 1.211$\times$; because no maximum A/D gate was registered, it is
reported as dispersion rather than converted into a pass/fail criterion.

\begin{table}[t]
  \caption{Prepared public \rtdl versus direct CUDA/OptiX (A/D). Ratio is
  lower-is-better; times are medians of worker medians.}
  \label{tab:prepared}
  \small
  \begin{tabular}{@{}llrrr@{}}
    \toprule
    GPU & Task & A (ms) & D (ms) & Median / max \\
    \midrule
    RTX 4090 & Relation & 0.2810 & 0.2612 & 1.0769 / 1.0923 \\
    RTX 4090 & Triangle & 0.0598 & 0.0510 & 1.1751 / 1.2110 \\
    RTX 3090 & Relation & 0.2402 & 0.2198 & 1.0948 / 1.1188 \\
    RTX 3090 & Triangle & 0.0608 & 0.0536 & 1.1336 / 1.1427 \\
    \bottomrule
  \end{tabular}
\end{table}

Across 20 observations, AOT cache-hit medians were 58.3--78.2~ms, about
1.04--2.02\% of cold time. Hits still include lookup, validation, and load and
are outside the prepared endpoint.

\subsection{Adverse lifecycle results}

Every A/C post-import median is adverse, 1.560--1.837$\times$, with a
2.377$\times$ worst block. Because entry was revised after the adverse result,
neither endpoint supports a confirmatory claim. Against E, A also regresses by
about 8--22\% at entry and 16--31\% post-import despite favorable prepared A/E,
showing material pre-steady-state protocol and Python/native lifecycle work.

\begin{table*}[t]
  \caption{Lifecycle ratios. A/C compares current \rtdl with the strong
  device-continuation baseline; A/E compares current with frozen predecessor
  control E. A/E entry and post-import rows are post hoc, non-gating
  diagnostics; only A/E prepared was registered. Ratios are lower-is-better.}
  \label{tab:lifecycle}
  \small
  \begin{tabular}{@{}llrrrrrr@{}}
    \toprule
    GPU & Task & A/C entry & A/C post median & A/C post range & A/E entry & A/E post & A/E prepared \\
    \midrule
    RTX 4090 & Relation & 0.654 & 1.749 & 1.725--1.866 & 1.192 & 1.305 & 0.584 \\
    RTX 4090 & Triangle & 0.642 & 1.560 & 1.527--1.639 & 1.080 & 1.169 & 0.903 \\
    RTX 3090 & Relation & 0.681 & 1.837 & 1.816--2.377 & 1.217 & 1.262 & 0.608 \\
    RTX 3090 & Triangle & 0.618 & 1.637 & 1.608--1.653 & 1.138 & 1.163 & 0.922 \\
    \bottomrule
  \end{tabular}
\end{table*}

Arm C outperformed B in all four formal steady comparisons (C/B ratios
0.221--0.654), establishing that B alone would be a weak baseline for this
experiment. It does not establish C as globally optimal. We also recorded
1,024 timer-instrumentation endpoints for A only. Paired overhead was measured
only for A; no corresponding paired overhead measurement was made for B, C, D,
or E.

No wrong output was observed in the final GPU samples. This statement is not a
correctness proof and does not repair the detailed-receipt shortfall described
in Section~\ref{sec:lifecycle}.

\subsection{Application-route cost against public PyOptiX}
\label{sec:app-pyoptix}

An initial frozen $A_0$/PyOptiX transaction on this RTX A4500 covered selected
stages from three of nine mappings: Particle, RT-2A1, and LibRTS point/range.
All twelve complete, first-execute, and prepared ratios were adverse:
1.080--10.208$\times$ complete, 2.732--75.533$\times$ first-execute, and
1.379--39.083$\times$ prepared. We preserve this result rather than replacing
it with the successor.

Remediation removed repeated compilation, materialization, and receipt work and
reused identity-bound executable families. A new transaction was frozen before
worker zero. Successor RTDL ($A_s$) and PyOptiX ($C$) both exclude offline
compilation. Complete includes deployment through close; prepared reuses state
after one untimed warmup. Each row has eight paired fresh-process blocks, four
per order; the envelope is median $\leq1.20$ and every block $\leq1.35$.

\begin{table*}[t]
  \caption{Remediated RTX A4500 routes. Each endpoint is $A_s/C$ ms; paired
  ratio [block min--max]; paired $A_s/A_0$. Lower is better; $A_0$ is immutable
  old V4. All ten rows meet the registered envelope.}
  \label{tab:app-pyoptix-results}
  \scriptsize
  \centering
  \begin{tabular}{@{}lrrl@{}}
    \toprule
    Unit & Complete & Prepared & Output \\
    \midrule
    Particle & 449.2/2401.6; .185 [.112--.201]; .019 & .402/.542; .773 [.684--.840]; .013 & $5000\!\times\!3$ U32 \\
    com-dblp/1M & 3764.2/3636.3; 1.055 [.952--1.165]; .328 & 478.4/435.1; 1.100 [1.050--1.183]; .807 & 2,224,385 \\
    LibRTS point & 1448.3/4433.5; .323 [.292--.357]; .375 & .256/.251; .999 [.733--1.110]; .024 & 112,729 \\
    LibRTS range & 1629.4/4451.7; .365 [.331--.420]; .405 & .252/.259; .985 [.612--1.243]; .027 & 105,826 \\
    cit-Patents/4M & 12712.2/12257.5; 1.038 [1.003--1.058]; .455 & 9347.9/8802.2; 1.057 [1.038--1.084]; .867 & 7,515,023 \\
    \bottomrule
  \end{tabular}
\end{table*}

All 240 workers had distinct PIDs, equal outputs, CPU-8 affinity, and zero
retry/discard/timeout; an independent recount reconstructed all 80 cells. The
preselected cit-Patents row executes 28 segments and 69,850,749 physical
queries: prepared $A_s/C$ is 9.348/8.802~s and 1.057$\times$; other rows are
subsecond. Complete starts from derived inputs. Particle is a fixed app-shaped
specialization; Graph and LibRTS use typed families but retain app-owned
encoding and oracles. The result covers five units and one host condition, not
the six unmeasured mappings or an app-independent engine.

\section{Limitations and Broader Implications}

Users provide bounded callback intent and data; shared \rtdl machinery validates
schemas, cross-role effects, ABI, identity, and lifecycle; trusted per-topology
lowerers generate and bind CUDA/OptiX code. Failed status gates suppress output,
while experiments may add application oracles.

This split does not make an arbitrary application correct. The author must
select an RT formulation, encode geometry, declare trusted semantic meanings,
provide domain predicates and reductions not covered by a fixed route, and
validate the result against an independent oracle. The compiler can reject
inconsistency among represented obligations, but a wrong yet internally
coherent declaration can pass. Nor does verified Callback IR imply an
executable path: only topologies with audited lowerers can be materialized.

The language trades expressiveness for analysis. It excludes dynamic
allocation, recursion, reflection, unrestricted loops, raw pointers, and
arbitrary foreign calls. Its stable facade has two constructors. The
two-generation synthetic paths use exact-IR specializations; the application
study also measures disclosed callback-leaf and closed native routes. No
independent developer has measured authoring effort. These are design costs,
not temporary omissions that the evaluation silently treats as solved.

The implementation suggests, but does not validate elsewhere, a pattern for
managed-language/accelerator boundaries: use the application-visible result
contract to organize callback legality, target-action interpretation, and
publication; then distinguish metadata checks from trusted lowering obligations
and execution-time failures. Other accelerators may expose analogous splits,
but their effects, lowerers, and checker vocabularies would need independent
design and evidence.

\section{Related Work}

\subsection{Repurposed RT applications and abstractions}

A scoped recent review reports 35 proposed RT solutions across 32 problems and
already identifies coherent design of geometry, queries, and hit-produced
operations as a recurring concern~\cite{meneses2026rtcoresurvey}. RayJoin,
RayDB, RTNN, RT-DBSCAN, graph triangle counting, X-HD, particle tracking, and
RT-BarnesHut each supply application-specific mappings and result handling
~\cite{geng2024rayjoin,shi2025raydb,zhu2022rtnn,nagarajan2023rtdbscan,xiao2025rtgraph,geng2026xhd,wang2022particle,nagarajan2025rtbarneshut}.
They establish both the opportunity and prior solutions to individual
obligations; RTDL does not claim to discover geometric encoding, all-hit
continuation, deduplication, nearest-K maintenance, or graph orientation.

LibRTS is the closest application-facing abstraction in this set. It provides
point and range spatial queries plus user device handlers for counting and
collection~\cite{geng2025librts}. RTDL is narrower in supported operations and
broader in a different dimension: it accepts restricted role source, represents
typed effects and result-route constraints, and binds the accepted program to
generated target and runtime identities. We have not run a head-to-head user
study, and do not claim superiority or that LibRTS cannot add such checks.

TTA/TTA+ instead changes RT hardware and its interface to support more general
tree traversal~\cite{ha2024tta}. RTDL works with existing OptiX hardware after
the author has selected an RT mapping; it neither removes fixed-function limits
nor finds a profitable mapping automatically.

\subsection{RT languages and compilation systems}

OptiX established programmable RT, including early termination, continued
intersection processing, attributes, and non-rendering examples
~\cite{parker2010optix}. Current OptiX, Vulkan, and DXR specify substantial
pipeline, interface, payload, and SBT rules
~\cite{nvidiaOptix9Guide,khronosVulkanInterfaces,khronosVulkanRT,microsoftDXRPAQ}.
PyOptiX exposes the OptiX host API to Python, while OWL owns much pipeline
construction~\cite{nvidiaPyOptix,nvidiaOWL}. These are not unchecked baselines.

Shader Components and Slang provide modular interfaces, specialization, and
reflection; Slang capabilities validate target, stage, API, and hardware
requirements, and shader cursors address parameter layout and binding
~\cite{he2017shaderComponents,he2018slang,slangCapabilities,slangShaderCursors}.
SlangPy adds Python RT pipelines, hit groups, shader tables, dispatch,
marshalling, caching, and lifecycle support
~\cite{slangPyRelease,slangPyDocs}. Luisa automates dependency analysis and
RT pipeline/SBT/launch work; CrossRT generates host/device and RT mappings from
restricted C++; Dr.Jit traces and specializes broad Python/C++ computation
through OptiX
~\cite{zheng2022luisa,luisaCompute,frolov2024crossrt,jakob2022drjit}.
Scion is a DSL/compiler that separates BVH logical trees, physical layouts, and
traversal algorithms with static well-formedness checks
~\cite{gyurgyik2026scion}. These systems preclude claims that RTDL first
automates RT construction, relates logic to physical structure, or specializes
code from result use.

Rendering DSLs also encode domain contracts. Open Shading Language separates
shader-produced closures from renderer-owned integration, and shader-component
systems coordinate code with required parameters
~\cite{openShadingLanguageDocs,he2017shaderComponents}. RTDL applies a bounded
language design to selected non-rendering result protocols; it does not claim
that rendering lacks cross-stage semantics.

The concrete RTDL increment is its implemented connection, for supported fixed
families, from declared result obligations to role-effect admission, semantic
and physical contracts, trusted traversal-action generation, executable
identity, and interface-specific failure checks. It uses explicit family rules
and guarded specialization at the cost of restricted expressiveness and
topology-specific trusted code. The available sources for adjacent systems do
not establish this exact combination, but source silence is not an inability
result and those systems may be extensible.

\subsection{Protocol and validation foundations}

Typestate, interface automata, and multiparty session types provide general and
often formally stronger models of state and global/local compatibility
~\cite{strom1986typestate,deAlfaro2001interfaceAutomata,honda2016multiparty}.
Typed linking and FFI checking already relate cross-language behavior,
representations, offsets, tags, and effects
~\cite{furr2005ffi,patterson2017linkingTypes}. Proof-carrying code established
consumer-side policy validation of supplied proofs before native execution
~\cite{necula1996safe}. RTDL contributes neither a new linking calculus nor a
proof certificate. Schemas, projections, lowerers, and runtime remain trusted;
hashes establish identity, not semantics. Its finite target checker is
empirical and structural, unlike formal GPU verification such as
GPUVerify~\cite{betts2012gpuverify}.

\section{Artifact}

The nine-member evidence-only archive contains an anonymous projection and
standard-library verifier over 160 formal cells, 20,480 steady samples, 1,024
A-only instrumentation endpoints, 20 AOT samples, and eight competence records.
Two clean exports were byte-identical and replayed in normal and optimized
Python from fresh directories. It bundles no third-party code, raw GPU archive,
or CUDA/OptiX material and performs offline recount, not GPU rerun,
installation, or private-history reconstruction.

Application transactions are separate. The private 1,049-file successor
archive supports clean-checkout recount of raw hashes, identities, outputs,
and statistics. A separate seven-file anonymous projection retains all 240
worker rows, 1,248 timed samples, 120 warmups, and 80 paired cells; its
standard-library verifier reconstructs the ten evaluations. The projection
removes private path, host, GPU-UUID, process-ID, and source identities. Neither
package reruns GPU work or includes every raw input, and the anonymous
projection does not replace private custody.

\section{Threats to Validity}

\paragraph{Construct validity.}
Prepared A/D excludes authoring, cold compilation, and first use. Direct D is a
named contract-equivalent implementation, not a proven latency lower bound;
strong C is output-equivalent but not operation-identical. These observations
do not isolate intrinsic language overhead, and paired instrumentation covers A
only. Application complete starts after loaders; the successor excludes offline
compilation from both arms, but routes and timed validation still differ.
The same-width semantic-ABI example is illustrative: mutation tests establish
check liveness, but no retained run executes its defective route. A trusted
semantic schema may also be wrong while internally coherent.

\paragraph{Internal validity.}
Both synthetic tasks were adapted, not unseen. Entry was revised after an
adverse result; both first-result endpoints remain lifecycle-confounded and are
not reinterpreted as passes. Mocks show that a provider double fault can replace
the primary exception and lose retry ownership, while a native fork can evade
the cached-PID guard. Neither mock ran GPU work or appeared in retained workers.
CPU 8 followed a disclosed diagnosis; an unlocked run failed two short-row
gates. An unpooled 48-CPU scan put Particle CPU 8 at 1.400 versus formal .773,
so sub-ms rows are unstable; GPU clocks were unlocked.

\paragraph{External validity.}
Synthetic measurements cover two NVIDIA generations and two exact tasks without
cross-machine raw-time comparison. The remediated application study covers
selected stages from 3/9 apps, four original units, one long graph scale, and
one RTX A4500. Six apps are unexecuted; inputs are derived; peak memory was not
captured. There is no independent-user, prevalence, or arbitrary-Python/IR
evidence. Six recovered entry points also lack retained historical per-file
identities and the original archive; older timings are not pooled with M.

\paragraph{Implementation trust.}
The sphere route required about 2,635 topology-specific TCB lines; the finite
checker missed an early-return probe; 4,096 timed A calls lack per-call receipts.
The application remediation adds trusted executable-family and admission code;
Particle still relies on an app-shaped fixed standard-library native route.

\paragraph{Artifact scope.}
The nine-member F2 projection and separate seven-file application projection
check their disclosed identities, counts, formulas, and rows, but cannot
recreate GPU runs or private custody. The latter cross-binds absent private
bytes by hash rather than independently validating their contents.

\section{Conclusion}

\rtdl is a restricted-Python DSL and compiler for bounded repurposed ray
tracing. Authors express typed data and role behavior while retaining ownership
of the application mapping and semantic oracle. Typed Callback IR and
whole-protocol admission connect selected result obligations to role effects,
semantic and physical contracts, continuation, and executable identity.
Generated Numba leaves and trusted topology wrappers implement supported
authored templates; closed routes can select dedicated native implementations,
including exact-IR specializations for the measured standard routes.

The design produces concrete compiler behavior: it rejects a role-legal effect
that conflicts with one complete-relation route, generates physical
intersection ignore for logically accepted all-hit events, and exits a fixed
specialization when source semantics change. Nine historical V4 mappings show
where these components compose and where application responsibility remains;
they do not prove arbitrary generalization or usability. Across two exact tasks
and two GPU generations, prepared median latency is
1.077--1.175$\times$ Direct, while post-import comparison reaches
1.837$\times$ median and a 2.377$\times$ worst block. An initial selected-stage
application transaction was adverse at all twelve endpoints. Its separately
identified remediation put all ten registered complete/prepared rows within a
1.20$\times$-median/1.35$\times$-maximum envelope; the multi-second graph row
was 1.057$\times$ public PyOptiX. This establishes practical cost for the tested
bounded families and preserves substantial coverage, trust, and authoring
limits; it is not broad performance or productivity evidence.

\bibliographystyle{ACM-Reference-Format}
\bibliography{references}

\end{document}
